% ============================================================================
%  SECCIÓN 2.2: DESARROLLO DE LOS MARCOS DEL PROYECTO
% ============================================================================

\section{Desarrollo de los marcos del proyecto}

Los marcos del proyecto se desarrollan a continuaci\'on siguiendo una l\'ogica
deductiva, es decir, partiendo de los fundamentos generales de la ingenier\'ia
de software y de los sistemas de informaci\'on geogr\'afica hasta llegar a los
aspectos particulares del sistema de visualizaci\'on interactiva de escenarios
deportivos. Cada marco aporta el sustento te\'orico, conceptual, hist\'orico y
legal de los elementos que se presentan en el cap\'itulo I y en el cap\'itulo
III, de modo que las decisiones de dise\~no y desarrollo se apoyen en
conocimiento cient\'ificamente validado \citep{arias2012, hernandez2014}.

\subsection{Marco te\'orico}

El marco te\'orico integra las teor\'ias, los estudios y los enfoques
te\'oricos de validez que permiten el correcto encuadre del tema de
investigaci\'on. Se organiza en antecedentes de la investigaci\'on, rese\~na de
la instituci\'on, bases te\'oricas y sistema de variables.

\subsubsection{Antecedentes de la investigaci\'on}

Los antecedentes de la investigaci\'on consisten en la presentaci\'on de la
informaci\'on m\'as relevante y directamente relacionada con el tema del
proyecto, cuyo estudio garantiza que el trabajo parta del punto m\'as actual de
la tem\'atica en cuesti\'on \citep{arias2012}. Su revisi\'on, desarrollada en la
secci\'on de antecedentes del cap\'itulo I, evidencia que la informaci\'on sobre
la infraestructura deportiva del pa\'is se encuentra fragmentada y que no
existe un registro unificado, georreferenciado y de acceso p\'ublico que
permita conocer la oferta deportiva nacional.

Los antecedentes hist\'oricos muestran que la infraestructura deportiva del
pa\'is es considerable y de larga tradici\'on, pero su difusi\'on ha dependido
hist\'oricamente de los medios de cada recinto. Los antecedentes tem\'aticos
demuestran, mediante plataformas internacionales, que la centralizaci\'on de la
informaci\'acion georreferenciada incrementa la visibilidad de los recintos y
favorece el uso de la infraestructura existente \citep{longley2015,
pressman2010}. Por su parte, las investigaciones previas en el \'ambito de los
sistemas de informaci\'on geogr\'afica (SIG) aplicados a servicios urbanos han
confirmado que la visualizaci\'on cartogr\'afica de la infraestructura mejora la
accesibilidad a la informaci\'on y apoya la planificaci\'on de pol\'iticas
p\'ublicas \citep{longley2015}.

Estos antecedentes fundamentan directamente el problema planteado en la
secci\'on de planteamiento del problema del cap\'itulo I y la necesidad de
desarrollar el sistema que se propone en el cap\'itulo III, confirmando que la
soluci\'on no parte de cero, sino que se apoya en experiencias nacionales e
internacionales documentadas.

\subsubsection{Rese\~na de la instituci\'on}

El proyecto se desarrolla en el marco de la asignatura Aplicaciones M\'oviles
I de la carrera de Ingenier\'ia en Sistemas, Facultad de Ingenier\'ia, de la
Universidad Privada Domingo Savio (UPDS), sede Cochabamba. La UPDS es una
instituci\'on de educaci\'on superior privada que forma profesionales
competentes en las \'areas de la ingenier\'ia y las tecnolog\'ias de la
informaci\'on, vinculando el proceso formativo con el desarrollo de soluciones
reales para la sociedad.

En este contexto, el proyecto constituye un trabajo de investigaci\'on
aplicada que integra los contenidos de la asignatura \textemdash{} desarrollo
m\'ovil multiplataforma, interfaces de usuario y despliegue de
aplicaciones\textemdash{} con las necesidades del \'ambito deportivo nacional.
La instituci\'on brinda el marco acad\'emico, el acompa\~namiento docente y el
entorno de evaluaci\'on en el que se valida el desarrollo del sistema,
conforme a la estructura establecida en la gu\'ia de proyectos de grado.

\subsubsection{Bases te\'oricas}

Las bases te\'oricas constituyen el cuerpo de conocimiento cient\'ifico que
sustenta el proyecto. A continuaci\'on se desarrollan los conceptos te\'oricos
m\'as relevantes, organizados de lo general a lo particular y citando las
fuentes especializadas en las que se apoyan.

\paragraph{Sistemas de informaci\'on.} Un sistema de informaci\'on es un
conjunto de componentes interrelacionados que recolectan, procesan, almacenan
y distribuyen informaci\'on para apoyar la toma de decisiones y el control de
una organizaci\'on \citep{kendall2011}. Los sistemas de informaci\'on pueden
ser manuales o computarizados, y sus componentes b\'asicos son las personas,
los datos, los procedimientos, el hardware y el software \citep{kendall2011}.

La ingenier\'ia de software, disciplina que se ocupa del desarrollo de estos
sistemas de forma sistem\'atica, define el proceso de an\'alisis, dise\~no,
codificaci\'on y prueba de productos de software con calidad y dentro de
costos y plazos razonables \citep{pressman2010, sommerville2011}. La
importancia de estos fundamentos radica en que la plataforma propuesta es un
sistema de informaci\'on orientado a la difusi\'on de la infraestructura
deportiva: recibe, procesa y entrega informaci\'on sobre escenarios a
los usuarios, por lo que su construcci\'on debe seguir las buenas pr\'acticas
de la ingenier\'ia de software y del an\'alisis de sistemas \citep{kendall2011,
pressman2010}. Estas bases sustentan las secciones de metodolog\'ia y de
identificaci\'on de usuarios del cap\'itulo I y III respectivamente.

\paragraph{Ingenier\'ia de requerimientos.} La ingenier\'ia de requerimientos
es el proceso de establecer los servicios que el sistema debe proporcionar y
las restricciones bajo las cuales debe operar \citep{sommerville2011}. Los
requerimientos se clasifican en funcionales, que describen lo que el sistema
debe hacer, y no funcionales, que describen las propiedades del sistema como
el rendimiento, la seguridad y la usabilidad \citep{sommerville2011,
kendall2011}. Un proceso correcto de identificaci\'on de requerimientos reduce
los errores de desarrollo y asegura que el producto satisfaga las necesidades
reales de los usuarios \citep{sommerville2011}.

En el presente proyecto, la ingenier\'ia de requerimientos se aplic\'o para
definir los requerimientos funcionales y no funcionales preliminares que se
presentan en el cap\'itulo III, y se fundamenta en las t\'ecnicas de
recolecci\'on de datos declaradas en el cap\'itulo I (encuesta, entrevista y
revisi\'on documental), tal como lo establece el objetivo espec\'ifico de
identificaci\'on de requerimientos.

\paragraph{Arquitectura de software cliente--servidor.} La arquitectura
cliente--servidor es un modelo de computaci\'on distribuida en el que los
clientes (aplicaciones m\'oviles o navegadores web) solicitan recursos y
servicios a un servidor central, que procesa las peticiones y devuelve los
resultados \citep{sommerville2011, kendall2011}. Este modelo permite separar
la interfaz de usuario de la l\'ogica de negocio y del almacenamiento de
datos, favoreciendo la escalabilidad, la seguridad y el mantenimiento del
sistema \citep{sommerville2011}.

La arquitectura del sistema propuesto sigue este modelo: la aplicaci\'on
m\'ovil constituye el cliente, el backend (Node.js con
Express) act\'ua como servidor de la l\'ogica de negocio y el motor de base de
datos (PostgreSQL con la extensi\'on espacial PostGIS) gestiona la
persistencia de la informaci\'on de los escenarios deportivos. Esta base
te\'orica sustenta la selecci\'on de tecnolog\'ias preliminares presentada en el
cap\'itulo III.

\paragraph{Desarrollo \'agil e iterativo.} El desarrollo \'agil es un enfoque
que prioriza la entrega temprana y continua de funcionalidades de valor,
adapt\'andose a los cambios de manera iterativa e incremental \citep{larman2004,
sommerville2011}. Entre sus principios se encuentran la colaboraci\'on con el
cliente, la entrega frecuente de software funcional y la mejora continua del
proceso \citep{sommerville2011}. El desarrollo iterativo permite reducir el
riesgo y obtener retroalimentaci\'on temprana de los usuarios \citep{larman2004}.

En el marco del proyecto, el enfoque \'agil se materializa en la definici\'on
del producto m\'inimo viable (MVP), presentada en el cap\'itulo III, que
contempla \'unicamente las funcionalidades indispensables para que la soluci\'on
funcione y pueda ser probada por los usuarios. Este enfoque sustenta adem\'as
la justificaci\'on metodol\'ogica del cap\'itulo I, que propone validar la
propuesta con usuarios reales en etapas tempranas \citep{larman2004}.

\paragraph{Usabilidad y experiencia de usuario.} La usabilidad es el atributo
de calidad que mide la facilidad con la que los usuarios utilizan una
interfaz \citep{nielsen1994}. Seg\'un Nielsen, la usabilidad se compone de
cinco atributos: facilidad de aprendizaje, eficiencia de uso, facilidad de
recordaci\'on, tolerancia a errores y satisfacci\'on subjetiva
\citep{nielsen1994}. Una interfaz usable debe dise\~narse centrada en el
usuario, es decir, considerando las necesidades, capacidades y limitaciones de
quienes la utilizar\'an \citep{nielsen1994}.

La usabilidad es un elemento central del presente proyecto: la interfaz de
la aplicaci\'on m\'ovil debe ser intuitiva para asistentes, gestores y
administradores sin experiencia previa, tal como se establece en el
requerimiento no funcional de usabilidad del cap\'itulo III y en la variable
de visualizaci\'on interactiva definida en el cap\'itulo I. Las pruebas de
usabilidad previstas en las t\'ecnicas de investigaci\'on se basan en los
principios de esta base te\'orica.

\paragraph{Calidad de software.} La calidad del software se define a trav\'es
de modelos est\'andar que caracterizan los atributos que debe reunir un
producto de software. El modelo ISO/IEC 25010 establece caracter\'isticas de
calidad como la funcionalidad, la eficiencia, la compatibilidad, la
usabilidad, la fiabilidad, la seguridad, la mantenibilidad y la portabilidad
\citep{iso25010}. La aplicaci\'on de estos modelos permite evaluar de manera
objetiva el grado en que el sistema cumple con los requerimientos
\citep{iso25010}.

En el proyecto, el modelo ISO/IEC 25010 sirve de referencia para los
requerimientos no funcionales del cap\'itulo III (rendimiento, seguridad,
compatibilidad, portabilidad y calidad de software) y para la evaluaci\'on del
sistema propuesta en el objetivo espec\'ifico de evaluaci\'on de usabilidad,
rendimiento y satisfacci\'on del cap\'itulo I.

\paragraph{Sistemas de informaci\'on geogr\'afica.} Un sistema de informaci\'on
geogr\'afica (SIG) es un conjunto de herramientas dise\~nadas para capturar,
almacenar, manipular, analizar, gestionar y presentar informaci\'on
geogr\'aficamente referenciada \citep{longley2015}. Los SIG integran la base de
datos con capacidades de an\'alisis espacial y representaci\'on cartogr\'afica,
permitiendo relacionar datos de distinta naturaleza a trav\'es de su posici\'on
geogr\'afica \citep{longley2015}. Su aplicaci\'on en servicios urbanos,
turismo y planificaci\'on ha demostrado mejoras significativas en la
accesibilidad a la informaci\'on y en la toma de decisiones \citep{longley2015}.

Esta base te\'orica es fundamental para el proyecto: el sistema debe ubicar
geogr\'aficamente los escenarios deportivos y representarlos en un mapa, tal
como se define en la variable de visualizaci\'on interactiva del cap\'itulo I y
en los requerimientos funcionales de mapa interactivo y cat\'alogo del
cap\'itulo III. Los conceptos de georreferenciaci\'on y de datos espaciales se
desarrollan a continuaci\'on.

\paragraph{Georreferenciaci\'on y datos espaciales.} La georreferenciaci\'on
consiste en asociar a un elemento una posici\'on \'unica en la superficie
terrestre, expresada normalmente mediante coordenadas de latitud y longitud
\citep{longley2015}. Los datos geogr\'aficos se representan en un SIG mediante
modelos vectoriales, que codifican los elementos como puntos, l\'ineas y
pol\'igonos con sus atributos asociados \citep{longley2015}. Para almacenar y
consultar eficientemente estos datos se utilizan bases de datos espaciales,
que incorporan tipos de datos y funciones geogr\'aficas especializadas, como la
extensi\'on PostGIS de PostgreSQL, que permite calcular distancias y \'areas y
determinar la cercan\'ia de un punto a otros elementos \citep{postgresql}.

En el proyecto, cada escenario deportivo se georreferencia mediante sus
coordenadas, lo que permite al usuario consultar los escenarios m\'as cercanos
a su ubicaci\'on, tal como se requiere en la funcionalidad de mapa interactivo
del cap\'itulo III y en el indicador de cobertura del sistema definido en el
cap\'itulo I.

\paragraph{Visualizaci\'on interactiva de datos geoespaciales.} La
visualizaci\'on de datos es el proceso de representar informaci\'on de forma
gr\'afica para facilitar su comprensi\'on y an\'alisis \citep{longley2015}. En el
\'ambito geoespacial, la visualizaci\'on interactiva permite al usuario
explorar un mapa, aplicar zoom, desplazarse y consultar la informaci\'on
asociada a cada elemento representado \citep{longley2015}. Las librer\'ias de
mapas interactivos, como Leaflet, proporcionan los componentes necesarios
(marcadores, capas, popups y controles de navegaci\'on) para construir este
tipo de visualizaciones de forma ligera y compatible con m\'ultiples
plataformas \citep{leaflet}.

La visualizaci\'on interactiva constituye el n\'ucleo de la propuesta: permite
representar los escenarios deportivos sobre un mapa, de modo que el usuario
identifique r\'apidamente los recintos de su inter\'es y acceda a su informaci\'on
detallada. Esta base sustenta la variable de visualizaci\'on interactiva del
cap\'itulo I, el objetivo general del proyecto y las funcionalidades de mapa
interactivo y detalle del escenario del cap\'itulo III.

\paragraph{Desarrollo m\'ovil multiplataforma.} El desarrollo de aplicaciones
m\'oviles ha evolucionado desde las aplicaciones nativas, desarrolladas con el
lenguaje y las herramientas propias de cada plataforma, hasta los enfoques
multiplataforma, que permiten construir una \'unica base de c\'odigo ejecutable
en m\'ultiples sistemas operativos, reduciendo costos y tiempos de desarrollo
\citep{pressman2010}. React Native es un framework de c\'odigo abierto que
permite desarrollar aplicaciones nativas para iOS y Android utilizando
JavaScript y el modelo de componentes de React \citep{reactnative}. Expo es
una plataforma construida sobre React Native que simplifica la configuraci\'on,
la compilaci\'on y la publicaci\'on de la aplicaci\'on, adem\'as de facilitar su
ejecuci\'on en dispositivos f\'isicos mediante la aplicaci\'on Expo Go
\citep{expo}.

Esta base te\'orica sustenta la selecci\'on de React Native y Expo como
tecnolog\'ias del frontend m\'ovil en el cap\'itulo III y el requerimiento no
funcional de compatibilidad (funcionamiento en iOS y Android) definido en el
mismo cap\'itulo.

\paragraph{Backend y servicios en la nube.} El backend de una aplicaci\'on
m\'ovil concentra la l\'ogica de negocio y la comunicaci\'on con la base de
datos. Node.js es un entorno de ejecuci\'on de JavaScript orientado a
aplicaciones de red con alto rendimiento y escalabilidad \citep{nodejs}. Los
servicios en la nube ofrecen infraestructura y plataformas gestionadas que
permiten desplegar aplicaciones sin administrar servidores propios; Supabase,
por ejemplo, proporciona autenticaci\'on, base de datos y almacenamiento
integrados con un modelo de servicio gratuito adecuado para proyectos de
investigaci\'on \citep{supabase}.

En el proyecto, estas tecnolog\'ias sustentan la arquitectura cliente--servidor
definida en el cap\'itulo III, garantizando la disponibilidad del sistema y el
manejo seguro de los datos de los usuarios, tal como lo exigen los
requerimientos no funcionales de seguridad y disponibilidad.

\paragraph{Control de versiones.} El control de versiones es una pr\'actica de
la ingenier\'ia de software que permite registrar, organizar y rastrear los
cambios realizados sobre el c\'odigo fuente de un proyecto \citep{gitdocs}. Git
es el sistema de control de versiones distribuido m\'as utilizado, y GitHub es
la plataforma de alojamiento remoto que facilita la colaboraci\'on entre los
integrantes del equipo y la trazabilidad de los cambios \citep{githubdocs}.

Esta base te\'orica sustenta la configuraci\'on de Git y GitHub del cap\'itulo
III, que garantiza la colaboraci\'on ordenada del equipo, el respaldo del
c\'odigo y el seguimiento del desarrollo del sistema a lo largo de las etapas
del proyecto.

\subsubsection{Sistema de variables}

El sistema de variables relaciona los conceptos te\'oricos con los aspectos
medibles de la realidad. Una variable es un aspecto de la realidad que cambia
en distintas situaciones y para diferentes objetivos y sujetos; las variables
permiten agrupar, diferenciar, ordenar, distribuir y relacionar los elementos
de la realidad \citep{hernandez2014}.

Para el presente proyecto se identificaron tres variables, cuya
operacionalizaci\'on completa se presenta en la secci\'on de definici\'on de
variables del cap\'itulo I: acceso a la informaci\'on deportiva, visualizaci\'on
interactiva y cobertura del sistema. La variable de acceso a la informaci\'on
deportiva se sustenta en las bases te\'oricas de sistemas de informaci\'on y de
difusi\'on; la variable de visualizaci\'on interactiva se sustenta en las bases
te\'oricas de SIG, georreferenciaci\'on y visualizaci\'on interactiva; y la
variable de cobertura del sistema se sustenta en las bases te\'oricas de bases
de datos espaciales y arquitectura de software \citep{hernandez2014}.

\subsection{Marco conceptual}

El marco conceptual asigna un significado preciso y contextualizado a los
conceptos, t\'erminos y variables principales del tema de investigaci\'on
\citep{arias2012}. Su prop\'osito es uniformizar el lenguaje utilizado en el
documento y delimitar el alcance de cada t\'ermino. En la tabla se presentan
los conceptos principales utilizados en el proyecto, con su definici\'on y la
fuente que los sustenta.

\begin{table}[H]
  \centering
  \caption{Marco conceptual del proyecto}
  \contenidoConFuente{%
    \begin{tablaAPA}
    \begin{tabular}{@{}p{4.0cm}p{9.0cm}p{2.8cm}@{}}
      \toprule
      \textbf{Concepto} & \textbf{Definici\'on} & \textbf{Fuente} \\
      \midrule
      Sistema de informaci\'on &
      Conjunto de componentes interrelacionados que recolectan, procesan,
      almacenan y distribuyen informaci\'on para apoyar la toma de decisiones &
      \citet{kendall2011} \\
      \addlinespace
      Ingenier\'ia de software &
      Disciplina que aplica principios de ingenier\'ia al desarrollo
      sistem\'atico de software con calidad &
      \citet{pressman2010} \\
      \addlinespace
      Sistema de informaci\'on geogr\'afica &
      Herramientas para capturar, almacenar, analizar y presentar informaci\'on
      geogr\'aficamente referenciada &
      \citet{longley2015} \\
      \addlinespace
      Georreferenciaci\'on &
      Asociaci\'on de un elemento a una posici\'on \'unica en la superficie
      terrestre mediante coordenadas &
      \citet{longley2015} \\
      \addlinespace
      Visualizaci\'on interactiva &
      Representaci\'on gr\'afica de datos que permite al usuario explorarlos y
      modificarlos de forma din\'amica &
      \citet{longley2015} \\
      \addlinespace
      Mapa interactivo &
      Mapa digital que permite al usuario navegar, hacer zoom y consultar la
      informaci\'on asociada a cada elemento &
      \citet{leaflet} \\
      \addlinespace
      Escenario deportivo &
      Recinto o infraestructura destinada a la pr\'actica de una o m\'as
      disciplinas deportivas &
      Elaboraci\'on propia \\
      \addlinespace
      Aplicaci\'on multiplataforma &
      Aplicaci\'on que se ejecuta en m\'ultiples sistemas operativos a partir de
      un \'unico c\'odigo base &
      \citet{pressman2010} \\
      \addlinespace
      Producto m\'inimo viable (MVP) &
      Versi\'on m\'as peque\~na de un producto que aporta valor real y permite
      validar la soluci\'on con retroalimentaci\'on temprana &
      \citet{larman2004} \\
      \addlinespace
      Requerimiento funcional &
      Servicio que el sistema debe proporcionar a sus usuarios &
      \citet{sommerville2011} \\
      \addlinespace
      Requerimiento no funcional &
      Restricci\'on o propiedad que el sistema debe cumplir, como rendimiento,
      seguridad o usabilidad &
      \citet{sommerville2011} \\
      \addlinespace
      Usabilidad &
      Facilidad con la que los usuarios aprenden y utilizan una interfaz &
      \citet{nielsen1994} \\
      \addlinespace
      Calidad de software &
      Grado en que un producto satisface los atributos de calidad establecidos
      en modelos est\'andar &
      \citet{iso25010} \\
      \addlinespace
      Wireframe &
      Esquema de baja fidelidad que representa la estructura y los elementos de
      una interfaz &
      \citet{pressman2010} \\
      \addlinespace
      Cliente--servidor &
      Modelo de computaci\'on distribuida en el que los clientes solicitan
      servicios a un servidor central &
      \citet{sommerville2011} \\
      \bottomrule
    \end{tabular}
    \end{tablaAPA}%
  }{Elaboraci\'on propia en base a la revisi\'on documental.}
\end{table}

Los conceptos presentados se utilizan de manera consistente en el cap\'itulo I
y en el cap\'itulo III del presente informe, garantizando la coherencia
terminol\'ogica del documento.

\subsection{Marco hist\'orico}

El marco hist\'orico recopila los antecedentes hist\'oricos referentes al tema
de estudio, cit\'andolos en orden cronol\'ogico \citep{arias2012}. Su
construcci\'on permite comprender la evoluci\'on de la infraestructura deportiva
y de las tecnolog\'ias de informaci\'on que sustentan la propuesta.

\paragraph{La infraestructura deportiva del pa\'is.} La pr\'actica deportiva
organizada en el pa\'is se consolid\'o durante la primera mitad del siglo XX,
con la construcci\'on de los primeros estadios y coliseos en las principales
ciudades, impulsada por el desarrollo del f\'utbol y de otras disciplinas. Con
el paso de las d\'ecadas, la oferta de recintos se diversific\'o hasta incluir
estadios, coliseos cerrados, canchas de baloncesto y voleibol, frontones,
piscinas, pistas de atletismo y complejos polideportivos, distribuidos en los
nueve departamentos del pa\'is, como se detalla en los antecedentes hist\'oricos
del cap\'itulo I.

\paragraph{La evoluci\'on de los sistemas de informaci\'on.} La segunda mitad
del siglo XX estuvo marcada por la aparici\'on de las computadoras y por el
desarrollo de la ingenier\'ia de software como disciplina, que evolucion\'o
desde enfoques artesanales hacia procesos sistem\'aticos y orientados a la
calidad \citep{pressman2010}. En la d\'ecada de 1980 surgieron los sistemas de
informaci\'on geogr\'afica, que integraron la informaci\'on tabular con la
representaci\'on cartogr\'afica, abriendo nuevas posibilidades para el an\'alisis
espacial \citep{longley2015}.

\paragraph{La era de internet y los sistemas distribuidos.} La masificaci\'on
de internet en la d\'ecada de 1990 impuls\'o el desarrollo de sistemas web y de
la arquitectura cliente--servidor, permitiendo que la informaci\'on se
distribuyera a escala global \citep{sommerville2011}. En el \'ambito deportivo,
esta etapa correspondi\'o a la difusi\'on de la informaci\'on mediante sitios web
y, posteriormente, mediante redes sociales, aunque de forma fragmentada y sin
una visi\'on integral de la infraestructura existente.

\paragraph{La era m\'ovil y los mapas interactivos.} A partir de la d\'ecada de
2010, la masificaci\'on de los tel\'efonos inteligentes y de los mapas digitales
interactivos transform\'o la manera en que las personas acceden a la
informaci\'on geogr\'afica \citep{longley2015}. El desarrollo multiplataforma,
representado por frameworks como React Native y plataformas como Expo,
permiti\'o construir aplicaciones para iOS y Android con un \'unico c\'odigo base
\citep{reactnative, expo}. En el \'ambito nacional, sin embargo, la
informaci\'on sobre la infraestructura deportiva sigui\'o dispersa y
desactualizada, sin aprovechar estas tecnolog\'ias.

Este recorrido hist\'orico muestra que las condiciones tecnol\'ogicas y
sociales actuales hacen viable y pertinente la propuesta de una aplicaci\'on
m\'ovil de visualizaci\'on interactiva de escenarios deportivos, tal como se
plantea en los objetivos del cap\'itulo I y en la propuesta del cap\'itulo III.

\subsection{Marco referencial}

El marco referencial recopila las normas legales relacionadas con el tema de
investigaci\'on \citep{arias2012}. El proyecto se enmarca en el ordenamiento
jur\'idico del Estado Plurinacional de Bolivia, que reconoce el deporte como un
derecho y regula el uso de las tecnolog\'ias de la informaci\'on y la
comunicaci\'on. Las principales disposiciones aplicables son las siguientes:

\paragraph{Constituci\'on Pol\'itica del Estado Plurinacional de Bolivia.} La
Constituci\'on, en su art\'iculo sobre los derechos de la salud y el deporte,
reconoce el derecho de toda persona al deporte y a la pr\'actica de la
actividad f\'isica, y establece la obligaci\'on del Estado de promover y
fomentar la cultura deportiva y el acceso de la poblaci\'on a la infraestructura
deportiva \citep{cpe2009}. Este mandato constitucional fundamenta la
finalidad social del proyecto: facilitar el acceso de la ciudadan\'ia a la
informaci\'on sobre los escenarios deportivos para promover la pr\'actica
deportiva, en l\'inea con la justificaci\'on social del cap\'itulo I.

\paragraph{Ley N.\textordmasculine{} 277 de 2012, Ley del Deporte.} La Ley
del Deporte establece las pol\'iticas, los principios y las normas para la
organizaci\'on y la promoci\'on del deporte en el pa\'is, as\'i como las
obligaciones de los \'organos competentes en materia de construcci\'on,
mantenimiento y administraci\'on de escenarios deportivos \citep{ley277}. La
ley reconoce la importancia de la infraestructura deportiva para el desarrollo
del deporte y promueve la participaci\'on de las entidades p\'ublicas y privadas
en su gesti\'on. Esta norma respalda la pertinencia de visibilizar y difundir
la oferta de escenarios deportivos del pa\'is, tal como lo hace el sistema
propuesto.

\paragraph{Ley N.\textordmasculine{} 164 de 2011, Ley General de
Telecomunicaciones, Tecnolog\'ias de Informaci\'on y Comunicaci\'on.} La Ley
N.\textordmasculine{} 164 regula el uso de las tecnolog\'ias de informaci\'on y
comunicaci\'on (TIC) en el pa\'is y reconoce su importancia para el desarrollo
econ\'omico y social \citep{ley164}. La ley promueve el acceso a las TIC y el
desarrollo de servicios y aplicaciones digitales que beneficien a la
poblaci\'on. Esta norma sustenta la base legal del desarrollo de la aplicaci\'on
m\'ovil propuesta, as\'i como las decisiones tecnol\'ogicas presentadas en el
cap\'itulo III.

En conjunto, el marco referencial proporciona el sustento normativo del
sistema, que debe respetar la normativa vigente en materia de deporte y de uso
de tecnolog\'ias de la informaci\'on, garantizando que la soluci\'on propuesta se
encuentre alineada con el ordenamiento jur\'idico del pa\'is.
